\chapter{Modèles de taux forward et de marché}\label{ch:rc:forward-rate-and-market-models}

Pendant vingt ans avant 1997, les traders évaluaient les caps avec la formule de Black appliquée
à chaque taux forward comme s'il s'agissait d'une action, sachant que la formule provenait
d'aucun modèle cohérent de la courbe~: le forward lognormal d'un caplet et le
forward lognormal du suivant ne pouvaient être tous deux lognormaux sous la même
mesure, et les \omterm{def:rc:short-rate-models:short}{modèles de taux court} évaluaient les caps avec d'autres formules. En 1997, trois
articles, de Brace, Gatarek et Musiela, de Jamshidian, et de Miltersen,
Sandmann et Sondermann ont montré que la formule des traders était exacte dans
un modèle dans lequel chaque taux forward est lognormal sous la mesure de sa propre
date de paiement, et qu'un seul tel modèle pouvait tenir tous les forwards à la fois. Le
modèle de marché, comme on l'appela ensuite, se calibre sur les caplets par construction,
demande Monte Carlo pour presque tout le reste, et rend visible un paramètre que les
\omterm{def:rc:short-rate-models:short}{modèles de taux court} masquaient~: la corrélation entre les taux forward. Ce chapitre
dérive le cadre général auquel il appartient, le construit et le simule, et montre
deux corrélations qui ajustent les mêmes caplets et évaluent la même swaption 22~\%
d'écart.

\section{Le cadre de Heath--Jarrow--Morton et son drift}

\begin{definition}[Cadre de Heath--Jarrow--Morton]\label{def:rc:forward-rate-and-market-models:hjm}
Le \emph{cadre de Heath--Jarrow--Morton}\index{cadre de Heath--Jarrow--Morton}
(1992) modélise toute la courbe des \omterm{def:rc:curve-construction:pillar}{taux forward instantanés}~: pour chaque échéance $T$,
$df(t,T) = \alpha(t,T)\,dt+\sigma_f(t,T)\,dW_t$ sous $\mathbb Q$, à partir de
la courbe d'aujourd'hui $f(0,T)$, avec des volatilités $\sigma_f(t,T)$ choisies librement (vecteurs
pour plusieurs facteurs).
\end{definition}

\begin{definition}[Condition de drift HJM]\label{def:rc:forward-rate-and-market-models:drift}
La \emph{condition de drift HJM}\index{condition de drift HJM} est la restriction
qu'impose l'absence d'arbitrage sur le drift une fois les volatilités
choisies~:
\[
\alpha(t,T) = \sigma_f(t,T)\int_t^T\sigma_f(t,u)\,du .
\]
\end{definition}

\begin{proposition}[Pourquoi le drift est fixé]\label{prop:rc:forward-rate-and-market-models:drift}
Les prix des obligations $P(t,T) = \exp(-\int_t^Tf(t,u)\,du)$ actualisés par le compte bancaire
sont des $\mathbb Q$-martingales si et seulement si la condition de drift est vérifiée.
\end{proposition}
\begin{proof}
$d\ln P(t,T) = f(t,t)\,dt-\int_t^Td f(t,u)\,du = \bigl(r_t-\int_t^T\alpha(t,u)\,du\bigr)dt
-\Sigma(t,T)\,dW_t$ avec $\Sigma(t,T)=\int_t^T\sigma_f(t,u)\,du$. Par la formule d'Itô
le drift de $P/B$ vaut $-\int_t^T\alpha\,du+\frac12\Sigma^2$~; il s'annule pour tout
$T$ si et seulement si $\int_t^T\alpha(t,u)\,du = \frac12\Sigma(t,T)^2$, et en dérivant par rapport à $T$
on obtient la condition.
\end{proof}

\begin{example}[Hull--White est un modèle HJM]\label{ex:rc:forward-rate-and-market-models:hw}
Le modèle du chapitre~7 a $\sigma_f(t,T) = \sigma e^{-\kappa(T-t)}$~: son drift vaut
$\sigma^2e^{-\kappa(T-t)}(1-e^{-\kappa(T-t)})/\kappa$, soit 4.1 points de base par an pour le
forward à dix ans avec $\sigma = 80$ points de base et $\kappa = 3~\%$. Les volatilités de
cette forme exponentielle sont exactement celles pour lesquelles la courbe forward est pilotée par
un état markovien unique~; une $\sigma_f$ générale rend le taux court dépendant du chemin
et le modèle ne peut plus qu'être simulé.
\end{example}

\section{Le modèle de marché}

Le \omterm{def:rc:curve-construction:pillar}{taux forward instantané} n'est pas négocié. Le modèle de marché retient les forwards
qui le sont~: $F_k(t)$ pour les périodes de calcul $[T_k,T_{k+1}]$ d'une structure de tenor.

\begin{definition}[Modèle de marché LIBOR]\label{def:rc:forward-rate-and-market-models:lmm}
Le \emph{modèle de marché LIBOR}\index{modèle de marché LIBOR} (LMM, ou modèle de marché forward
lognormal) rend chaque taux forward lognormal, avec une volatilité déterministe
sous la probabilité forward de sa date de paiement~:
$dF_k(t) = \sigma_k(t)F_k(t)\,dW^{T_{k+1}}_k(t)$, avec $d\langle W_i,W_j\rangle =
\rho_{ij}\,dt$. Son nom vient de l'indice interbancaire pour lequel il a été conçu~; la
dynamique vaut pour toute structure de tenors de taux prospectifs, et le chapitre~10
étend celle-ci à la capitalisation des taux au jour le jour.
\end{definition}

Chaque caplet s'évalue alors exactement par la formule de Black, avec une volatilité égale à
$\sqrt{\frac1{T_k}\int_0^{T_k}\sigma_k(t)^2dt}$, ce qui fait précisément que le modèle
se calibre sur les caps par construction. Une forme courante de volatilité est donnée par la \omterm{def:rc:forward-rate-and-market-models:rebonato}{formule de Rebonato}
\[
\sigma_k(t) = \phi_k\,g(T_k-t),\qquad g(\tau) = (a+b\tau)e^{-c\tau}+d ,
\]
une bosse en fonction du temps avant fixation, multipliée par une constante propre à chaque forward
(la \cref{fig:rc:forward-rate-and-market-models:abcd})~; les $\phi_k$ sont déterminés de sorte
que chaque caplet soit de nouveau évalué à son prix de marché.

\begin{omfigure}
\begin{tikzpicture}
\begin{axis}[width=11cm,height=5cm,
  xlabel={temps avant la fixation $\tau$ (années)},ylabel={$g(\tau)$},
  tick label style={font=\small},label style={font=\small},
  xmin=0,xmax=10,ymin=0,ymax=0.35,grid=major,grid style={gray!25},
  table/col sep=comma]
\addplot[omThm,very thick] table[x=tau,y=g]{figdata/rates-credit-risk/08-forward-rate-and-market-models/abcd.csv};
\end{axis}
\end{tikzpicture}

\omcaption{Fonction de volatilité de Rebonato avec $a=0.05$, $b=0.10$, $c=0.60$,
$d=0.15$~: la volatilité d'un forward est maximale environ un an et demi avant sa
fixation. Ajustée forward par forward, elle fait dépendre la volatilité de chaque forward de
son temps avant fixation, si bien que la courbe de volatilité observée dans un an ressemble à celle d'aujourd'hui.
Données~: le tutoriel du chapitre.}\label{fig:rc:forward-rate-and-market-models:abcd}
\end{omfigure}

\section{Drifts sous les mesures spot et terminale}

Pour simuler tous les forwards ensemble, ils doivent partager une même mesure, et sous toute
mesure unique, tous sauf un d'entre eux acquièrent un drift.

\begin{definition}[Mesure spot, mesure terminale]\label{def:rc:forward-rate-and-market-models:measures}
La \emph{mesure spot}\index{mesure spot} $\mathbb Q^d$ d'une structure de tenors a pour
numéraire le compte bancaire rééquilibré de façon discrète $B_d(t) =
P(t,T_{m(t)})\prod_{j<m(t)}(1+\delta_jF_j(T_j))$, qui détient l'obligation venant à échéance à
la prochaine date de fixation $T_{m(t)}$ et se reporte sur la suivante à chaque fixation. La
\emph{mesure terminale}\index{mesure terminale} est la probabilité forward de la dernière
date de paiement $T_n$, avec pour numéraire $P(t,T_n)$.
\end{definition}

\begin{proposition}[Dérives des modèles de marché]\label{prop:rc:forward-rate-and-market-models:spotdrift}
Sous la \omterm{def:rc:forward-rate-and-market-models:measures}{mesure spot}, pour un forward $F_k$ encore vivant en $t$,
\[
\frac{dF_k}{F_k} = \sigma_k(t)\sum_{j=m(t)}^{k}\frac{\delta_j\rho_{kj}\sigma_j(t)F_j}{1+\delta_jF_j}\,dt+\sigma_k(t)\,dW_k,
\]
et, sous la \omterm{def:rc:forward-rate-and-market-models:measures}{mesure terminale}, la dérive vaut $-\sigma_k\sum_{j=k+1}^{n-1}
\delta_j\rho_{kj}\sigma_jF_j/(1+\delta_jF_j)$.
\end{proposition}
\begin{proof}[Démonstration partielle]
On change de numéraire une période à la fois~: la densité de $\mathbb Q^{T_{j+1}}$ par
rapport à $\mathbb Q^{T_j}$ est proportionnelle à $P(t,T_{j+1})/P(t,T_j) =
1/(1+\delta_jF_j)$, dont la volatilité vaut $-\delta_j\sigma_jF_j/(1+\delta_jF_j)$~; par
le théorème de Girsanov (livre~4, chapitre~5), chaque changement ajoute à la dérive la
covariance de $\ln F_k$ avec cette densité. En sommant depuis l'obligation du numéraire
jusqu'à la date de paiement propre à $F_k$, on obtient les deux formules.
\end{proof}

\begin{omfigure}
\begin{tikzpicture}[x=1.25cm,font=\small]
\draw[->,thick] (0,0) -- (9.6,0) node[right] {temps};
\foreach \x in {0,...,8} {\draw (\x,-0.1) -- (\x,0.1); \node[below,font=\footnotesize] at (\x,-0.12) {$T_{\x}$};}
\foreach \x in {0,...,7} {\draw[omDef,fill=omDef!10] (\x+0.06,0.35) rectangle (\x+0.94,0.75); \node[font=\scriptsize] at (\x+0.5,0.55) {$F_{\x}$};}
\draw[omThm,very thick,-{Stealth[length=2mm]}] (0,1.25) -- (1,1.25) -- (1,1.55) -- (2,1.55) -- (2,1.85) -- (3,1.85);
\node[text=omThm,anchor=west,font=\footnotesize] at (3.1,1.85) {mesure spot~: on reporte l'obligation suivante à chaque fixation};
\draw[omProp,very thick,-{Stealth[length=2mm]}] (0,2.6) -- (8,2.6);
\node[text=omProp,anchor=south,font=\footnotesize] at (4,2.65) {mesure terminale~: l'obligation venant à échéance à la dernière date};
\end{tikzpicture}

\omcaption{Une structure de tenors de huit forwards annuels et les numéraires des
deux mesures de simulation usuelles. Sous la \omterm{def:rc:forward-rate-and-market-models:measures}{mesure spot}, la dérive de chaque forward
somme les forwards compris entre la fixation suivante et lui-même~; sous la \omterm{def:rc:forward-rate-and-market-models:measures}{mesure
terminale}, ceux qui lui succèdent.}\label{fig:rc:forward-rate-and-market-models:tenor}
\end{omfigure}

On simule $\ln F_k$ avec la dérive évaluée au début du pas (Euler), puis de nouveau
à la fin prédite, en prenant la moyenne des deux
(prédicteur--correcteur)~: la dérive dépend des forwards eux-mêmes, et la correction
maintient les caplets simulés sur leurs prix de Black lorsque les pas sont grossiers.

\omcode{code/firm/lmm/firm_lmm.py}{86}{114}{La dérive sous la mesure spot et le pas prédicteur--correcteur.}\label{lst:rc:forward-rate-and-market-models:step}

\begin{example}[Les caplets reviennent]\label{ex:rc:forward-rate-and-market-models:caplets}
Dix forwards annuels issus de la courbe OIS en euros du chapitre~2 (1.98~\% pour la
première année, puis 3.03~\% pour la dixième), \omterm{def:rc:vanilla-rates-options:flat}{volatilités de caplets} de 70 points de
base en volatilité normale converties en Black (33.0~\% à un an, puis 23.1~\% à neuf ans), fonction de
Rebonato avec les paramètres de la \cref{fig:rc:forward-rate-and-market-models:abcd}
et corrélation $\rho_{ij}=e^{-0.1|i-j|}$~: 20\,000 trajectoires sous la \omterm{def:rc:forward-rate-and-market-models:measures}{mesure spot}, quatre
pas par an, reproduisent les caplets à la monnaie à moins de deux erreurs types
(\cref{fig:rc:forward-rate-and-market-models:caplets}).
\end{example}

\begin{omfigure}
\begin{tikzpicture}
\begin{axis}[width=11cm,height=5.2cm,
  xlabel={échéance du caplet (années)},ylabel={volatilité de Black (\%)},
  tick label style={font=\small},label style={font=\small},
  xmin=0.5,xmax=9.5,ymin=20,ymax=36,grid=major,grid style={gray!25},
  legend columns=2,legend style={font=\footnotesize,at={(0.5,-0.25)},anchor=north,draw=none,fill=none,/tikz/every even column/.append style={column sep=8pt}},
  table/col sep=comma]
\addplot[omDef,thick] table[x=k,y=target]{figdata/rates-credit-risk/08-forward-rate-and-market-models/caplets.csv};
\addplot[omThm,only marks,mark=*,mark size=2pt,error bars/.cd,y dir=both,y explicit] table[x=k,y=mc,y error=err]{figdata/rates-credit-risk/08-forward-rate-and-market-models/caplets.csv};
\legend{volatilité de caplet cible,Monte Carlo ($\pm2$ erreurs types)}
\end{axis}
\end{tikzpicture}

\omcaption{\omterm{def:rc:vanilla-rates-options:flat}{Volatilités de caplets} à la monnaie implicites des prix Monte Carlo du
modèle de marché calibré, comparées aux cibles. Données~: le
tutoriel du chapitre.}\label{fig:rc:forward-rate-and-market-models:caplets}
\end{omfigure}

\section{Calibrer sur caps et swaptions~: la formule de Rebonato}

Un taux de swap est une moyenne pondérée de forwards,
$S_{a,b}=\sum_{i=a}^{b-1}w_iF_i$, avec les poids $w_i = \delta_iP(t,T_{i+1})/A_{a,b}(t)$~; ainsi
sa volatilité se déduit de celles des forwards et de leurs corrélations.

\begin{definition}[Modèle de marché des swaps]\label{def:rc:forward-rate-and-market-models:smm}
Un \emph{modèle de marché des swaps}\index{modèle de marché des swaps} rend lognormaux les taux de swap choisis
sous leurs propres mesures d'annuité, de sorte que les swaptions se prennent par
la formule de Black. Il ne peut pas coexister avec un modèle de marché lognormal des
forwards~: une somme pondérée de lognormales n'est pas lognormale. Les desks en retiennent un et
approchent l'autre.
\end{definition}

\begin{definition}[Formule de Rebonato]\label{def:rc:forward-rate-and-market-models:rebonato}
La \emph{formule de Rebonato}\index{formule de Rebonato} approche la volatilité de Black
de la swaption expirant en $T_a$ sur $S_{a,b}$ dans le modèle de marché, en figeant les
poids et les forwards à leurs valeurs d'aujourd'hui~:
\[
\sigma_{a,b}^2\,T_a \approx \frac1{S_{a,b}(0)^2}\sum_{i,j=a}^{b-1}
w_i(0)w_j(0)F_i(0)F_j(0)\,\rho_{ij}\int_0^{T_a}\sigma_i(t)\sigma_j(t)\,dt .
\]
\end{definition}

\omcode{code/firm/lmm/firm_lmm.py}{75}{83}{Volatilité de swaption à poids figés de Rebonato.}\label{lst:rc:forward-rate-and-market-models:rebonato}

\section{Corrélation et réduction factorielle}

\begin{definition}[Corrélation des taux forward]\label{def:rc:forward-rate-and-market-models:corr}
La \emph{corrélation des taux forward}\index{corrélation des taux forward} $\rho_{ij}$ d'un
modèle de marché est la corrélation instantanée des mouvements browniens qui pilotent
$F_i$ et $F_j$. Elle est en général paramétrée, par exemple
$\rho_{ij} = e^{-\beta|T_i-T_j|}$, puis réduite à deux ou trois facteurs en ne gardant que
les vecteurs propres dominants de la matrice (les composantes principales du
chapitre~3), ce qui accélère la simulation et élimine le bruit.
\end{definition}

Les caplets ne dépendent pas de la corrélation, les swaptions en dépendent, car un taux de swap moyenne
plusieurs forwards, et la moyenne de taux décorrélés en réduit la volatilité. Un modèle de marché
calibré sur les seuls caplets laisse donc le prix de ses swaptions à
une hypothèse.

\begin{example}[Mêmes caplets, swaptions différentes]\label{ex:rc:forward-rate-and-market-models:corr}
Calibrer le modèle sur les mêmes caplets avec $\beta=0.02$ (forwards presque
parfaitement corrélés) et avec $\beta=0.40$ (les forwards à un an et à neuf ans
corrélés à 0.04). La swaption d'échéance dans cinq ans sur un an est
pricée à 25.0~\% de volatilité de Black par les deux (c'est un seul forward)~; sur cinq ans,
à 22.1~\% et 17.3~\% (\cref{fig:rc:forward-rate-and-market-models:swaptions}). Sur
100~millions d'euros, à la monnaie (forward 2.933~\%, annuité 4.102), cela donne 2.35
millions d'euros contre 1.84 million~: un écart de 22~\% pour les mêmes caplets. La simulation de
Monte Carlo confirme la \omterm{def:rc:forward-rate-and-market-models:rebonato}{formule de Rebonato} à l'intérieur de ses erreurs types.
\end{example}

\begin{omfigure}
\begin{tikzpicture}
\begin{axis}[width=11cm,height=5cm,
  xlabel={indice forward $j$},ylabel={$\rho_{1j}$},
  tick label style={font=\small},label style={font=\small},
  xmin=1,xmax=9,ymin=0,ymax=1.05,grid=major,grid style={gray!25},
  legend columns=3,legend style={font=\footnotesize,at={(0.5,-0.25)},anchor=north,draw=none,fill=none,/tikz/every even column/.append style={column sep=8pt}},
  table/col sep=comma]
\addplot[omDef,thick,mark=*] table[x=j,y=b02]{figdata/rates-credit-risk/08-forward-rate-and-market-models/corr.csv};
\addplot[omProp,thick,mark=triangle*] table[x=j,y=b10]{figdata/rates-credit-risk/08-forward-rate-and-market-models/corr.csv};
\addplot[omThm,thick,mark=square*] table[x=j,y=b40]{figdata/rates-credit-risk/08-forward-rate-and-market-models/corr.csv};
\legend{$\beta=0.02$,$\beta=0.10$,$\beta=0.40$}
\end{axis}
\end{tikzpicture}

\omcaption{Corrélation du premier forward annuel avec les autres sous la
paramétrisation exponentielle $\rho_{ij}=e^{-\beta|T_i-T_j|}$ pour trois valeurs de
$\beta$. Données~: le tutoriel du chapitre.}\label{fig:rc:forward-rate-and-market-models:corrplot}
\end{omfigure}

\begin{omfigure}
\begin{tikzpicture}
\begin{axis}[width=11.5cm,height=5.4cm,
  xlabel={maturité du swap débutant dans cinq ans (années)},ylabel={volatilité de Black (\%)},
  tick label style={font=\small},label style={font=\small},
  xmin=0.8,xmax=5.2,ymin=15,ymax=27,xtick={1,2,3,4,5},grid=major,grid style={gray!25},
  legend columns=2,legend style={font=\footnotesize,at={(0.5,-0.25)},anchor=north,draw=none,fill=none,/tikz/every even column/.append style={column sep=8pt}},
  table/col sep=comma]
\addplot[omDef,thick] table[x=tenor,y=reb02]{figdata/rates-credit-risk/08-forward-rate-and-market-models/swaptions.csv};
\addplot[omDef,only marks,mark=*] table[x=tenor,y=mc02]{figdata/rates-credit-risk/08-forward-rate-and-market-models/swaptions.csv};
\addplot[omThm,thick,dashed] table[x=tenor,y=reb40]{figdata/rates-credit-risk/08-forward-rate-and-market-models/swaptions.csv};
\addplot[omThm,only marks,mark=square*] table[x=tenor,y=mc40]{figdata/rates-credit-risk/08-forward-rate-and-market-models/swaptions.csv};
\legend{Rebonato $\beta=0.02$,Monte Carlo $\beta=0.02$,Rebonato $\beta=0.40$,Monte Carlo $\beta=0.40$}
\end{axis}
\end{tikzpicture}

\omcaption{Volatilités de swaptions d'échéance cinq ans issues de deux modèles de marché qui
recalibrent les mêmes caplets. La décorrélation des forwards abaisse la volatilité des
taux de swap plus longs~; la \omterm{def:rc:forward-rate-and-market-models:rebonato}{formule de Rebonato} (traits) suit la simulation de Monte Carlo (marques).
Données~: le tutoriel du chapitre.}\label{fig:rc:forward-rate-and-market-models:swaptions}
\end{omfigure}

\begin{remark}[Calibrer la corrélation]\label{rem:rc:forward-rate-and-market-models:calib}
Les desks calibrent le modèle de marché conjointement~: les $\phi_k$ sur les caplets ou les
swaptions co-terminales, la forme de la volatilité et la corrélation sur la \omterm{def:rc:sabr-in-rates-and-the-volatility-cube:matrix}{matrice des
swaptions}, souvent en ajustant $\beta$ (ou une forme à deux paramètres) aux corrélations
historiques des taux forward et en laissant les swaptions fixer le reste. Un produit dépendant de la
corrélation des taux qu'aucun instrument liquide ne fixe porte une réserve
pour celle-ci (\cref{ch:rc:pnl-explain-and-independent-price-verification}).
\end{remark}

\section{Tutoriel~: un modèle de marché sur la courbe en euros}\label{tut:rc:forward-rate-and-market-models:tutorial}

\begin{tutorial}
\textbf{Objectif.} Calibrer un modèle de marché sur les caplets, le simuler et mesurer comment la
corrélation fait bouger les swaptions. \textbf{État final~:}
\cref{fig:rc:forward-rate-and-market-models:caplets,fig:rc:forward-rate-and-market-models:swaptions}.
\begin{tutsteps}
\item \textbf{Forwards et cibles}~: \texttt{F0} issu de la courbe €STR du chapitre~2, volatilités de Black
des caplets à partir de 70 points de base en normal.
\item \textbf{Calibrer}~: \texttt{calibrate\_to\_caplets(F0, vols, abcd, beta)}
résout les $\phi_k$.
\item \textbf{Simuler}~: \texttt{mc\_caplet} et \texttt{mc\_swaption} sous la \omterm{def:rc:forward-rate-and-market-models:measures}{mesure
spot}~; \texttt{caplet\_check()}.
\item \textbf{Corrélation}~: \texttt{swaption\_table()} pour $\beta=0.02$ et $0.40$~;
\texttt{fig\_rc\_lmm.py} écrit les graphiques.
\end{tutsteps}
\textbf{Pour aller plus loin.} Simuler sous la \omterm{def:rc:forward-rate-and-market-models:measures}{mesure terminale} et comparer les erreurs
types~; réduire la matrice de corrélation à deux facteurs et mesurer la variation de la
volatilité cinq dans cinq.
\end{tutorial}

\section{Construction~: un moteur de modèle de marché}\label{bld:rc:forward-rate-and-market-models:lmm}

\begin{build}
\bfield{Rôle} Le modèle de référence pour les produits dépendant de plusieurs taux
à plusieurs dates~: les bermudiennes (chapitre~9), les \omterm{def:rc:convexity-adjustments-and-constant-maturity-products:spread}{options sur spread CMS} (chapitre~6), et la
validation du modèle à un facteur du chapitre~7.
\bfield{Interface} \texttt{LMM(f0, phi, abcd, beta)} avec \texttt{vol},
\texttt{corr}, \texttt{integrated\_cov}, \texttt{discount\_factors}, \texttt{swap}~;
\texttt{calibrate\_to\_caplets}~; \texttt{rebonato\_vol}~; \texttt{simulate},
\texttt{mc\_caplet}, \texttt{mc\_swaption}~; \texttt{hjm\_drift\_gaussian}.
\bfield{Règles} Structure de tenor annuelle, forwards lognormaux, \omterm{def:rc:forward-rate-and-market-models:measures}{mesure spot} avec
étapes prédicteur--correcteur~; graines fixes~; numpy uniquement.
\bfield{Tests d'acceptation} \texttt{code/firm/lmm/tests/}~: les variances intégrées calibrées
égalent les cibles de caplets et la \omterm{def:rc:forward-rate-and-market-models:rebonato}{formule de Rebonato} pour une swaption à une période
renvoie la volatilité du caplet~; Monte Carlo reprice un caplet à moins de
trois erreurs types~; une swaption profondément dans la monnaie vaut le swap forward~; une
corrélation plus faible abaisse la volatilité des swaptions~; la dérive HJM de Hull--White
a sa forme fermée.
\bfield{Pour aller plus loin} Forwards décalés ou normaux pour les taux bas~; un modèle de marché
à volatilité stochastique~; calibration de $\beta$ sur la \omterm{def:rc:sabr-in-rates-and-the-volatility-cube:matrix}{matrice des swaptions}.
\end{build}

\begin{omsources}
\item D.\ Heath, R.\ Jarrow et A.\ Morton, ``Bond pricing and the term structure of
interest rates'', \emph{Econometrica} 60(1), 1992.
\item A.\ Brace, D.\ Gatarek et M.\ Musiela, ``The market model of interest rate
dynamics'', \emph{Mathematical Finance} 7, 1997 ; F.\ Jamshidian, ``LIBOR and swap
market models and measures'', \emph{Finance and Stochastics} 1, 1997 ; K.\
Miltersen, K.\ Sandmann et D.\ Sondermann, \emph{Journal of Finance} 52(1), 1997.
\item R.\ Rebonato, \emph{Modern Pricing of Interest-Rate Derivatives: The LIBOR
Market Model and Beyond}, Princeton University Press, 2002.
\end{omsources}

\section{Exercices}

\begin{exercise}[$\star$]\label{exo:rc:forward-rate-and-market-models:1}
Dans un modèle HJM à $\sigma_f(t,T)=\sigma$ constante, quelle est la dérive de $f(t,T)$~?
De quel \omterm{def:rc:short-rate-models:short}{modèle de taux court} s'agit-il~?
\end{exercise}

\begin{exercise}[$\star$]\label{exo:rc:forward-rate-and-market-models:2}
Pourquoi chaque caplet est-il évalué exactement par la formule de Black dans le modèle de marché,
mais pas chaque swaption~?
\end{exercise}

\begin{exercise}[$\star$]\label{exo:rc:forward-rate-and-market-models:3}
Avec $\beta = 0.40$, quelle est la corrélation entre le premier et le neuvième
forward annuel~? Et avec $\beta=0.02$~?
\end{exercise}

\begin{exercise}[$\star\star$]\label{exo:rc:forward-rate-and-market-models:4}
Sous la \omterm{def:rc:forward-rate-and-market-models:measures}{mesure spot}, écrire la dérive de $F_2$ en $t=0.5$ dans un modèle à
forwards annuels. Quels forwards interviennent~?
\end{exercise}

\begin{exercise}[$\star\star$]\label{exo:rc:forward-rate-and-market-models:5}
Deux forwards de 3~\% avec une volatilité de Black de 20~\% chacun, des pondérations de 0.5 chacune et
une corrélation $\rho$. Avec la \omterm{def:rc:forward-rate-and-market-models:rebonato}{formule de Rebonato}, quelle est la volatilité du taux de swap pour
$\rho=1$ et pour $\rho=0.5$~?
\end{exercise}

\begin{exercise}[$\star\star$]\label{exo:rc:forward-rate-and-market-models:6}
Lire la \cref{fig:rc:forward-rate-and-market-models:swaptions}~: pourquoi les deux modèles
s'accordent-ils sur la swaption cinq-into-un~?
\end{exercise}

\begin{exercise}[$\star\star\star$]\label{exo:rc:forward-rate-and-market-models:7}
\emph{Programmation.} Évaluer la swaption payeuse cinq-into-cinq à la monnaie sur 100 millions d'euros sous
les deux corrélations avec la \omterm{def:rc:forward-rate-and-market-models:rebonato}{formule de Rebonato}, puis donner la différence.
\end{exercise}

\begin{exercise}[$\star\star\star$]\label{exo:rc:forward-rate-and-market-models:8}
\emph{Trouver la faille.} «~Nous simulons le modèle de marché sous la \omterm{def:rc:forward-rate-and-market-models:measures}{mesure spot} en fixant la dérive
de chaque forward à zéro, car chaque forward est une martingale.~»
\end{exercise}

\section{Problème~: La corrélation que l'on ne voit pas}

\begin{problem}[{Problème du week-end --- un modèle de marché calibré sur les caplets}]\label{pb:rc:forward-rate-and-market-models:1}
Un desk fait tourner un modèle de marché annuel à dix forwards sur la courbe OIS en euros, calibré
chaque matin sur les caplets à 70 points de base de volatilité normale. Un client
demande la swaption payeuse cinq-into-cinq à la monnaie sur 100 millions d'euros. Le
paramètre de corrélation $\beta$ du desk n'a jamais été calibré.

\medskip\noindent\textbf{Partie I --- Ce que fixent les caplets.}
\begin{enumerate}
\item Donner le taux de swap forward et l'annuité.
\item Quels paramètres les caplets déterminent-ils, et lesquels ne déterminent-ils pas~?
\item Donner les volatilités de Black des caplets à un et à neuf ans.
\item Pourquoi la swaption cinq-into-un est-elle la même quelle que soit la corrélation~?
\item Que laisse encore libre la calibration~?
\end{enumerate}

\medskip\noindent\textbf{Partie II --- Deux corrélations.}
\begin{enumerate}[resume]
\item Donner la volatilité de la cinq-into-cinq avec $\beta = 0.02$ et avec $\beta=0.40$.
\item Donner les deux primes.
\item Donner l'écart relatif.
\item Dans quelle mesure la \omterm{def:rc:forward-rate-and-market-models:rebonato}{formule de Rebonato} concorde-t-elle ici avec Monte Carlo~?
\item Quelle corrélation attendrait-on des données historiques, et pourquoi~?
\end{enumerate}

\medskip\noindent\textbf{Partie III --- Y remédier.}
\begin{enumerate}[resume]
\item Quels instruments permettraient de fixer la corrélation~?
\item Si le marché cote la cinq-into-cinq à 20~\%, quel $\beta$ en découle-t-il
(l'encadrer)~?
\item Après avoir calibré $\beta$ sur ce cours, que resterait-il non fixé~?
\item Comment le choix de $\beta$ modifie-t-il les couvertures du desk~?
\item Pourquoi un \omterm{def:rc:short-rate-models:hw}{modèle de Hull--White} à un facteur équivaut-il à $\beta=0$~?
\end{enumerate}

\medskip\noindent\textbf{Partie IV --- Jugement.}
\begin{enumerate}[resume]
\item Quel prix faut-il coter, et quelle réserve faut-il constituer~?
\item Comment suivre ce paramètre dans le temps~?
\item Que doit dire la documentation du modèle sur $\beta$~?
\item Énoncer le \emph{résultat nommé}~: les primes de la cinq-into-cinq sous les deux
corrélations et leur écart.
\item En une phrase~: que sait une swaption que ne sait pas un caplet~?
\end{enumerate}
\end{problem}

\section{Questions d'entretien}

\begin{interviewq}[$\star$ \iqroles{chercheur}]\label{iq:rc:forward-rate-and-market-models:1}
Énoncez la \omterm{def:rc:forward-rate-and-market-models:drift}{condition de drift HJM} et expliquez d'où elle provient.
\end{interviewq}

\begin{interviewq}[$\star\star$ \iqroles{chercheur, développeur}]\label{iq:rc:forward-rate-and-market-models:2}
Comment simulez-vous un \omterm{def:rc:forward-rate-and-market-models:lmm}{modèle de marché LIBOR}, et pourquoi y a-t-il une dérive~?
\end{interviewq}

\begin{interviewq}[$\star\star$ \iqroles{chercheur, trader}]\label{iq:rc:forward-rate-and-market-models:3}
Pourquoi les caplets et les swaptions ne peuvent-ils pas être tous deux exactement lognormaux~?
\end{interviewq}

\begin{interviewq}[$\star\star$ \iqroles{trader}]\label{iq:rc:forward-rate-and-market-models:4}
Vous calibrez un modèle de marché sur les seuls caplets. Qu'est-ce qui cloche dans vos
prix de swaptions~?
\end{interviewq}

\begin{interviewq}[$\star\star\star$ \iqroles{chercheur, risque}]\label{iq:rc:forward-rate-and-market-models:5}
Comparez Hull--White et le modèle de marché pour évaluer un Bermuda à dix ans.
\end{interviewq}

\begin{interviewq}[$\star\star\star$ \iqroles{développeur}]\label{iq:rc:forward-rate-and-market-models:6}
Vos prix de caplets simulés s'écartent de Black lorsque vous utilisez des pas de temps
plus grossiers. Que se passe-t-il, et que modifiez-vous~?
\end{interviewq}
